\section{周期性}\label{節:周期関数の積分}

$T>0$ に対して $f(x+T)=f(x)$ がすべての $x$ で成り立つとき,\,$T$ を $f$ の周期という.周期という言葉だけから,\,$T$ は最小の正の周期だとは限らない.\,$\sin x$ の周期として $2\pi$ も $4\pi$ も使える.

周期関数をちょうど一周期分だけ積分した値は,始点によらない.まず $0\leq a\leq T$ の場合,
\[
\int_a^{a+T}f(x)\dd x
 =\int_a^T f(x)\dd x+\int_T^{a+T}f(x)\dd x.
\]
第二項で $x=t+T$ と置くと $f(t+T)=f(t)$ だから,
\[
\int_T^{a+T}f(x)\dd x=\int_0^a f(t)\dd t.
\]
二つを足すと
\begin{equation}\label{式:一周期の積分}
\int_a^{a+T}f(x)\dd x=\int_0^T f(x)\dd x
\end{equation}
になる.一般の $a$ も,整数個の $T$ だけ平行移動すれば始点を $[0,T)$ に移せるため,同じ結論になる.区間を長さ $T$ の部分に分ければ,正の整数 $n$ に対して
\[
\int_a^{a+nT}f(x)\dd x=n\int_0^T f(x)\dd x
\]
も得られる.

周期があるというだけでは,一周期分の積分が零とは限らない.たとえば $\sin^2x$ は周期 $\pi$ を持つが,一周期の積分は $\frac{\pi}{2}$ である.符号付きの部分が打ち消し合うか,積分区間が周期の整数倍になっているかを別々に確認する.

\begin{example}{絶対値がある周期関数}
次の定積分を求めよ.
\[
I=\int_{-\frac{\pi}{6}}^{\frac{13\pi}{6}}|\sin x|\dd x.
\]
\end{example}

$\sin(x+\pi)=-\sin x$ なので,絶対値を取ると周期は $\pi$ になる.一周期分は
\[
\int_0^\pi|\sin x|\dd x=\int_0^\pi\sin x\dd x=2.
\]
与えられた区間の長さは $\frac{7\pi}{3}$ であり,\,$2\pi$ と $\frac{\pi}{3}$ に分けられる.まず $-\frac{\pi}{6}$ から $\frac{11\pi}{6}$ までは二周期分なので積分は $4$ である.残りでは $x=t+2\pi$ と置き,
\[
\int_{\frac{11\pi}{6}}^{\frac{13\pi}{6}}|\sin x|\dd x
 =\int_{-\frac{\pi}{6}}^{\frac{\pi}{6}}|\sin t|\dd t.
\]
$|\sin t|$ は偶関数で,\,$0\leq t\leq\frac{\pi}{6}$ では $\sin t\geq0$ だから,
\[
I=4+2\int_0^{\frac{\pi}{6}}\sin t\dd t
 =4+2\left(1-\frac{\sqrt3}{2}\right)
 =6-\sqrt3.
\]
周期によって長い区間を短くし,偶奇によって符号の変わる区間を半分にした.周期の整数倍に満たない残りは,始点によって値が変わるので省いてはいけない.


